% =====================================================================
%  Manual del sistema --- Cancionero Infantil CENIDIM
%  Compilar con:  pdflatex Manual_CENIDIM.tex  (dos pasadas para el índice)
%  Solo requiere paquetes estándar de TeX Live / MiKTeX.
% =====================================================================
\documentclass[11pt,letterpaper,twoside]{report}

% ---------- Paquetes base ----------
\usepackage[utf8]{inputenc}
\usepackage[T1]{fontenc}
\usepackage[spanish,es-tabla,es-noindentfirst]{babel}
\usepackage[margin=2.6cm,headheight=14pt]{geometry}
\usepackage{xcolor}
\usepackage{graphicx}
\usepackage{booktabs}
\usepackage{longtable}
\usepackage{array}
\usepackage{enumitem}
\usepackage{fancyhdr}
\usepackage{titlesec}
\usepackage{listings}
\usepackage{microtype}
\usepackage{underscore}          % guiones bajos seguros y cortables en texto
\usepackage{amssymb}             % \checkmark, \square
\usepackage{lastpage}
\usepackage[colorlinks=true,
            linkcolor=cenidim-vino,
            urlcolor=cenidim-vino,
            citecolor=cenidim-vino]{hyperref}

% ---------- Paleta institucional (design tokens del frontend) ----------
\definecolor{cenidim-vino}{HTML}{751428}     % --color-brand
\definecolor{cenidim-vino-osc}{HTML}{5A0F1D} % --color-brand-dark
\definecolor{cenidim-mostaza}{HTML}{C5A46C}  % --color-brand-light
\definecolor{cenidim-crema}{HTML}{FAF7F0}    % --color-bg
\definecolor{cenidim-tinta}{HTML}{1A1612}    % --color-ink
\definecolor{cenidim-gris}{HTML}{8A7F6E}     % --color-text-muted

% ---------- Encabezado / pie ----------
\pagestyle{fancy}
\fancyhf{}
\fancyhead[LE]{\small\itshape Manual del sistema --- Cancionero CENIDIM}
\fancyhead[RO]{\small\itshape\leftmark}
\fancyfoot[C]{\small\thepage}
\renewcommand{\headrulewidth}{0.4pt}
\renewcommand{\headrule}{\color{cenidim-mostaza}\hrule width\headwidth height 0.4pt}

% ---------- Títulos ----------
\titleformat{\chapter}[display]
  {\normalfont\bfseries\color{cenidim-vino}}
  {\filleft\large\MakeUppercase{\chaptertitlename}\ \thechapter}
  {12pt}{\titlerule[1.2pt]\vspace{6pt}\huge}
\titlespacing*{\chapter}{0pt}{-20pt}{28pt}

% ---------- Listings ----------
\lstdefinestyle{shell}{
  basicstyle=\ttfamily\small,
  backgroundcolor=\color{cenidim-crema},
  frame=single, framerule=0.3pt, rulecolor=\color{cenidim-mostaza},
  breaklines=true, columns=fullflexible, keepspaces=true,
  xleftmargin=6pt, xrightmargin=6pt,
  aboveskip=10pt, belowskip=10pt,
}
\lstset{style=shell}

% ---------- Utilidades ----------
\providecommand{\nopagecolor}{\pagecolor{white}} % kernels antiguos
\newcommand{\ok}{\textcolor{cenidim-vino}{\textbf{\checkmark}}}
\newcommand{\pend}{\textcolor{cenidim-mostaza}{$\square$}}
\newcommand{\codigo}[1]{\texttt{#1}}
\emergencystretch=2em                     % evita overfull en texto justificado
\setlist{itemsep=2pt,topsep=4pt}

% =====================================================================
\begin{document}

% ==================== PORTADA ====================
\begin{titlepage}
  \pagecolor{cenidim-crema}
  \begin{center}
    \vspace*{1.5cm}

    {\color{cenidim-mostaza}\rule{\textwidth}{1.5pt}}\\[4pt]
    {\color{cenidim-vino}\rule{\textwidth}{0.6pt}}

    \vspace{2.2cm}
    {\small\scshape Centro Nacional de Investigación, Documentación\\[2pt]
     e Información Musical ``Carlos Chávez''}\\[1.8cm]

    {\Huge\bfseries\color{cenidim-vino} Cancionero Infantil\\[6pt] CENIDIM}\\[0.9cm]
    {\LARGE\color{cenidim-tinta} Manual del sistema}\\[0.5cm]
    {\large\itshape Plataforma web de consulta, análisis y administración\\ del archivo de letras de canciones}

    \vspace{2.6cm}
    \begin{tabular}{rl}
      \textbf{Versión del documento:} & 1.0 \\
      \textbf{Versión del sistema:}   & 0.2.0 \\
      \textbf{Fecha:}                 & Septiembre de 2026 \\
      \textbf{Rama de referencia:}    & \codigo{release} \\
      \textbf{Despliegue:}            & \codigo{cenidim.darylemb.dev} \\
    \end{tabular}

    \vfill
    {\color{cenidim-vino}\rule{\textwidth}{0.6pt}}\\[4pt]
    {\color{cenidim-mostaza}\rule{\textwidth}{1.5pt}}
  \end{center}
\end{titlepage}
\nopagecolor

% ---------- Portadilla ----------
\thispagestyle{empty}
\vspace*{6cm}
\begin{center}
  {\large\itshape Documento entregable del proyecto.}\\[4pt]
  {\small Describe la arquitectura, la instalación, el uso, la referencia
   de la API, el estado de calidad verificado y el trabajo pendiente.}
\end{center}
\cleardoublepage

% ==================== ÍNDICES ====================
\pagenumbering{roman}
\tableofcontents
\cleardoublepage
\listoftables
\cleardoublepage
\pagenumbering{arabic}

% =====================================================================
\chapter{Introducción}
% =====================================================================

\section{Propósito del sistema}
La plataforma \textbf{Cancionero CENIDIM} es una aplicación web que da
acceso público y administrado al archivo digital del cancionero
infantil del Centro Nacional de Investigación, Documentación e
Información Musical ``Carlos Chávez''. Permite:

\begin{itemize}
  \item Buscar canciones por título, álbum (fonograma) o contenido de
        la letra.
  \item Explorar el catálogo en una \emph{línea de tiempo} por año de
        registro.
  \item Consultar un \emph{tablero estadístico} con métricas del
        corpus: distribución por año, por tema, por clasificación
        lingüística, niveles de léxico (OOV) y nube de palabras.
  \item Administrar el catálogo (fonogramas y canciones), las cuentas
        de usuario y la auditoría operativa (roles
        \emph{viewer}/\emph{editor}/\emph{admin}).
\end{itemize}

\section{Cifras del corpus}
\begin{table}[h]
  \centering
  \begin{tabular}{lr}
    \toprule
    \textbf{Indicador} & \textbf{Valor} \\
    \midrule
    Canciones indexadas            & 3\,858 \\
    Fonogramas (álbumes)           & 280    \\
    Temas catalogados (distintos)  & $\sim$24 canónicos \\
    Clasificaciones lingüísticas   & 3 (estándar, regional, indígena) \\
    \bottomrule
  \end{tabular}
  \caption{Volumen del archivo al momento de esta edición.}
  \label{tab:corpus}
\end{table}

Los temas (\codigo{Tema:} en los archivos fuente) son asignados por
catalogadores humanos; el sistema \emph{no} los infiere. Las variantes
de escritura (mayúsculas, espacios, errores tipográficos conocidos) se
pliegan a formas canónicas mediante \codigo{canonical_tema}
(\codigo{backend/app/models/theme_normalization.py}).

\section{Alcance de este manual}
Este documento es el manual entregable del proyecto. Cubre la
arquitectura (cap.~\ref{cap:arquitectura}), el pipeline de datos
(cap.~\ref{cap:pipeline}), la instalación y el despliegue
(cap.~\ref{cap:instalacion}), el uso de la aplicación
(cap.~\ref{cap:usuario}), la referencia de la API
(cap.~\ref{cap:api}), el estado de calidad verificado
(cap.~\ref{cap:calidad}) y el trabajo pendiente
(cap.~\ref{cap:pendientes}).

% =====================================================================
\chapter{Arquitectura del sistema}
\label{cap:arquitectura}
% =====================================================================

\section{Vista general}
El sistema sigue una arquitectura de tres capas desplegada con
Docker Compose:

\begin{center}
\begin{tabular}{ccl}
  \toprule
  \textbf{Servicio} & \textbf{Puerto} & \textbf{Función} \\
  \midrule
  \codigo{frontend} & 80   & SPA Vue~3 servida por Nginx (sin root) \\
  \codigo{backend}  & 8000 & API FastAPI / uvicorn \\
  \codigo{db-init}  & ---  & Sidecar que regenera \codigo{letras.db} \\
  \bottomrule
\end{tabular}
\end{center}

El frontend consume la API bajo el prefijo \codigo{/api} (proxy Nginx en
producción, proxy Vite en desarrollo). La base de datos es
\textbf{SQLite} (\codigo{letras.db}), regenerada de forma determinista
por el sidecar \codigo{db-init} a partir de las fuentes, por lo que el
archivo binario es un artefacto de ejecución, no una fuente editada a
mano.

\section{Componentes}

\subsection{Backend --- FastAPI / Pydantic v2}
\begin{itemize}
  \item SQLAlchemy 2.0 (ORM) sobre el esquema \codigo{letras.db}
        (columnas \codigo{snake_case}); modelos en
        \codigo{app/models/}.
  \item Autenticación: hash de contraseñas sha256+bcrypt, JWT con
        \codigo{python-jose}, cookies HttpOnly + CSRF
        \emph{double-submit}, rotación de refresh tokens con revocación
        (tabla \codigo{RefreshTokenRevocation}).
  \item Autorización por roles:
        \codigo{viewer} $<$ \codigo{editor} $<$ \codigo{admin}, vía la
        fábrica \codigo{require_role(...)} en \codigo{app/deps.py}.
  \item Observabilidad: \codigo{/metrics} en formato Prometheus y
        logging JSON estructurado (\codigo{app/\allowbreak observability.py},
        \codigo{app/\allowbreak logging_config.py}).
  \item Migraciones con Alembic; en producción
        (\codigo{CENIDIM\_ENV=\allowbreak prod})
        el esquema \emph{no} se autocrea.
  \item Correo (recuperación de contraseña) vía Resend; en desarrollo,
        bandeja local en la tabla \codigo{email_outbox}.
\end{itemize}

\subsection{Frontend --- Vue 3 + TypeScript}
\begin{itemize}
  \item SPA con Vite, Vue Router 4, Pinia (stores \codigo{auth},
        \codigo{search}, \codigo{ui}, \codigo{filters}).
  \item Gráficas con \codigo{vue-chartjs} (Chart.js).
  \item TypeScript en modo \emph{strict}; estilos con design tokens
        (\codigo{src/assets/tokens.css}: Fraunces + Outfit + JetBrains
        Mono sobre paleta vino/mostaza/crema).
\end{itemize}

\subsection{Datos}
\begin{itemize}
  \item Fuentes: \codigo{db_fonografia.csv} (metadata de fonogramas) y
        \codigo{LetrasTXT/*.txt} (letras con línea \codigo{Tema:}).
  \item Artefacto: \codigo{letras.db} (SQLite) montado en
        \codigo{/data/letras.db} dentro del contenedor.
\end{itemize}

\section{Decisiones de arquitectura registradas}
El backend original en Go/Gin fue reemplazado por FastAPI y retirado
en la ``Fase~9'' del plan de migración. La decisión y sus
consecuencias se documentaron en el ADR
\codigo{docs/adr/\allowbreak 0001-fastapi-replaces-go.md} y el runbook
\codigo{docs/CUTOVER.md}; ambos archivos fueron retirados del árbol en
la limpieza posterior (commit \codigo{ff4303b}) y permanecen accesibles
en la historia de git. El backend de Go permanece
disponible en la historia de git (commit \codigo{2aab765}) solo como
referencia de emergencia.

% =====================================================================
\chapter{Pipeline de datos}
\label{cap:pipeline}
% =====================================================================

La base de datos se construye en tres pasos encadenados por
\codigo{scripts/build_db.sh} (o por el contenedor \codigo{db-init}):

\begin{enumerate}
  \item \textbf{Construcción} --- \codigo{scripts/build_db.py}: parsea
        \codigo{db_fonografia.csv} + \codigo{LetrasTXT/} y genera
        \codigo{letras.db} (port a Python del antiguo builder en Go,
        con salida byte-compatible). Crea el usuario administrador
        inicial si se define \codigo{ADMIN_PASS}.
  \item \textbf{Clasificación} --- \codigo{scripts/classify_songs.py}:
        con spaCy (\codigo{es_core_news_md}) calcula el porcentaje de
        palabras fuera de vocabulario (OOV) de cada letra y escribe
        \codigo{song_stats}:
        \begin{center}
        \begin{tabular}{lll}
          \toprule
          \textbf{Clasificación} & \textbf{Criterio} \\
          \midrule
          \codigo{ESPANOL_ESTANDAR}  & OOV $< 5\%$ \\
          \codigo{ESPANOL_REGIONAL}  & $5\% \le$ OOV $\le 18\%$ \\
          \codigo{LENGUA_INDIGENA}   & OOV $> 18\%$ o léxico indígena \\
          \bottomrule
        \end{tabular}
        \end{center}
  \item \textbf{Normalización} --- \codigo{scripts/normalize_db.py}:
        limpia el cuerpo de las letras (cabeceras, líneas
        \codigo{Dura:}/\codigo{Tema:}/\codigo{Personajes:}, iniciales),
        normaliza temas y revalida la correspondencia letra$\,\leftrightarrow\,$título.
        Es idempotente y conservadora (verificada con tests).
\end{enumerate}

\paragraph{Puntos de integridad verificados.}
Tras la normalización, 366 letras fueron limpiadas y 0 conservan
marcadores de metadatos residuales; el título y autor permanecen en
sus columnas dedicadas, por lo que la nube de palabras y el léxico OOV
reflejan únicamente el texto de la canción.

% =====================================================================
\chapter{Instalación y despliegue}
\label{cap:instalacion}
% =====================================================================

\section{Requisitos}
\begin{itemize}
  \item Docker + Docker Compose (vía recomendada), \emph{o}
  \item Python $\ge 3.12$ con \codigo{uv}; Node.js~24; spaCy
        \codigo{es_core_news_md} para regenerar la base de datos.
\end{itemize}

\section{Despliegue con Docker (recomendado)}
\begin{lstlisting}
docker compose up --build -d
backend/scripts/smoke.sh http://localhost:8000
\end{lstlisting}

Servicios resultantes:
\begin{itemize}
  \item Frontend: \codigo{http://localhost} (Nginx no privilegiado).
  \item API: \codigo{http://\allowbreak localhost:\allowbreak 8000/\allowbreak healthz};
        métricas en \codigo{/metrics}; especificación en
        \codigo{/\allowbreak openapi.json}.
\end{itemize}

Para producción (Coolify) se usa \codigo{docker-compose-coolify.yaml}
con volumen nombrado y secretos del gestor; el archivo por defecto
sigue siendo válido.

\section{Desarrollo local sin Docker}
\begin{lstlisting}
# Backend
cd backend
uv sync
PYTHONPATH=. uv run pytest tests/
uv run uvicorn app.main:app --port 8000 --reload

# Frontend (Node 24)
cd frontend
npm install
npm run dev          # Vite en :5173, proxy /api -> :8000
\end{lstlisting}

\section{Variables de entorno}
\begin{longtable}{p{5.2cm}p{2.1cm}p{7.2cm}}
  \toprule
  \textbf{Variable} & \textbf{Requerida} & \textbf{Descripción} \\
  \midrule
  \endhead
  \codigo{CENIDIM_JWT_SECRET} & Sí (prod) &
    Secreto de firma JWT ($\ge$ 256 bits;
    \codigo{openssl rand -hex 32}). Sin él el backend no arranca en
    producción. \\
  \codigo{CENIDIM_DB_PATH} & No & Ruta a \codigo{letras.db}
    (def. \codigo{letras.db}). \\
  \codigo{CENIDIM_ENV} & No & \codigo{dev} permite autocrear el
    esquema; \codigo{prod} exige Alembic. \\
  \codigo{CENIDIM_LOG_FORMAT} & No & \codigo{json} (producción) o
    \codigo{text}. \\
  \codigo{CENIDIM_WORKERS} & No & Workers de uvicorn (def. 2). \\
  \codigo{CORS_ALLOWED_ORIGINS} & No & Orígenes permitidos,
    separados por comas. \\
  \codigo{RESEND_API_KEY} & Para correo & Sin ella, los envíos caen en
    la tabla \codigo{email_outbox}. \\
  \codigo{CENIDIM_EMAIL_FROM} & Para correo & Remitente de correos de
    recuperación. \\
  \codigo{FRONTEND_BASE_URL} & Para correo & Origen del SPA para los
    enlaces de restablecimiento. \\
  \codigo{ADMIN_PASS} & Al sembrar & Contraseña del administrador
    inicial (build_db). \\
  \bottomrule
  \caption{Variables de entorno del backend.}
  \label{tab:envvars}
\end{longtable}

\section{Autenticación}
El inicio de sesión es usuario + contraseña (Google OAuth fue retirado
en la migración a FastAPI). \codigo{POST /api/auth/login} emite las
cookies HttpOnly \codigo{cenidim_session} (JWT de acceso),
\codigo{cenidim_refresh} y \codigo{cenidim_csrf}. Los endpoints
\codigo{/forgot} y \codigo{/reset} implementan recuperación por correo
(Resend o bandeja de desarrollo).

% =====================================================================
\chapter{Manual de usuario}
\label{cap:usuario}
% =====================================================================

\section{Vistas de la aplicación}
\begin{table}[h]
  \centering
  \begin{tabular}{lll}
    \toprule
    \textbf{Ruta} & \textbf{Vista} & \textbf{Acceso} \\
    \midrule
    \codigo{/}          & Línea de tiempo   & Público \\
    \codigo{/canciones} & Catálogo buscable & Público \\
    \codigo{/dashboards}& Tablero estadístico & Público \\
    \codigo{/admin}     & Panel de administración & JWT + rol \\
    \codigo{/login}     & Inicio de sesión  & Público \\
    \codigo{/forgot}, \codigo{/reset} & Recuperación & Público \\
    \bottomrule
  \end{tabular}
  \caption{Mapa de rutas del frontend.}
  \label{tab:rutas}
\end{table}

\subsection{Línea de tiempo}
Agrupa las canciones por año de registro del fonograma y permite abrir
cada pieza para leer su letra y metadatos.

\subsection{Catálogo (\emph{Canciones})}
Búsqueda paginada por título, álbum o contenido; filtros combinables
por tema, rango de años, clasificación lingüística y fonograma. La
columna de número no se trunca y el ordenamiento no produce parpadeo.

\subsection{Tablero estadístico}
KPIs del catálogo (canciones indexadas, álbumes en el filtro activo,
temas distintos), gráficas por año, por tema, por clasificación y por
nivel OOV, y nube de palabras por frecuencia. Cada gráfica tiene un
ícono \raisebox{0pt}{$\circledast$}\,(i) con explicación en lenguaje no
técnico; el texto fuente canónico es
\codigo{frontend/src/config/chartInfo.ts} (el glosario editorial
\codigo{docs/GLOSARIO.md} fue retirado del árbol en la limpieza de
\codigo{ff4303b}; disponible en la historia de git).

\subsection{Panel de administración}
CRUD de fonogramas, canciones y usuarios con paginación y ordenamiento
del lado del servidor; consulta de la bandeja de correo
(\codigo{email_outbox}) y de la bitácora de auditoría. La jerarquía de
roles es \codigo{viewer} (lectura) $<$ \codigo{editor} (escritura) $<$
\codigo{admin} (borrado y gestión de usuarios).

% =====================================================================
\chapter{Referencia de la API}
\label{cap:api}
% =====================================================================

La especificación completa y normativa es
\codigo{backend/openapi.json} (OpenAPI 3.1, regenerable con
\codigo{scripts/generate_openapi.py} y vigilada por un guardia de
deriva en CI). Resumen:

\section{Públicos}
\begin{longtable}{p{4.6cm}p{9.8cm}}
  \toprule
  \textbf{Endpoint} & \textbf{Descripción} \\
  \midrule
  \endhead
  \codigo{GET /healthz}     & Health check (Docker). \\
  \codigo{GET /metrics}     & Métricas Prometheus. \\
  \codigo{GET /openapi.json}& Especificación OpenAPI 3.1. \\
  \codigo{GET /api/search}  & Búsqueda paginada; \codigo{?query=} o \codigo{?q=}. \\
  \codigo{GET /api/song/\{id\}} & Detalle y letra de una canción. \\
  \codigo{GET /api/timeline}& Canciones agrupadas por año. \\
  \codigo{GET /api/stats}   & Métricas del tablero; honra los filtros
    compartidos (\codigo{tema}, \codigo{year_from}, \codigo{year_to},
    \codigo{clasificacion}, \codigo{album}, \codigo{q}). \\
  \codigo{GET /api/word-cloud} & Frecuencias de palabras (def. 200,
    máx. 500 vía \codigo{?limit=}). \\
  \bottomrule
  \caption{Endpoints públicos.}
  \label{tab:api-publica}
\end{longtable}

\section{Autenticación (\texttt{/api/auth/*})}
\begin{longtable}{p{4.6cm}p{9.8cm}}
  \toprule
  \textbf{Endpoint} & \textbf{Descripción} \\
  \midrule
  \endhead
  \codigo{POST /login}    & Emite cookies de sesión + \codigo{\{token, user\}}. \\
  \codigo{POST /register} & Auto-registro (si está habilitado). \\
  \codigo{POST /forgot}   & Siempre 200; envía enlace de un solo uso. \\
  \codigo{POST /reset}    & Consume el token y fija la nueva contraseña. \\
  \codigo{POST /refresh}  & Rota el refresh token (revoca el \codigo{jti} anterior). \\
  \codigo{POST /logout}   & Revoca el refresh token y limpia cookies. \\
  \codigo{GET /me}        & Perfil del usuario autenticado. \\
  \bottomrule
  \caption{Endpoints de autenticación.}
  \label{tab:api-auth}
\end{longtable}

\section{Administración (\texttt{/api/admin/*})}
Requieren JWT y rol suficiente. CRUD completo de \codigo{fonogramas},
\codigo{songs} y \codigo{users}; más los operativos
\codigo{GET /emails} (bandeja, filtro \codigo{only_failures}) y
\codigo{GET /audit} (bitácora, filtros \codigo{actor_id},
\codigo{action}).

% =====================================================================
\chapter{Calidad, pruebas e integración continua}
\label{cap:calidad}
% =====================================================================

\section{Estado verificado}
Los siguientes resultados fueron \emph{ejecutados y verificados} sobre
la rama \codigo{release} en la fecha de este documento:

\begin{table}[h]
  \centering
  \begin{tabular}{lll}
    \toprule
    \textbf{Verificación} & \textbf{Comando} & \textbf{Resultado} \\
    \midrule
    Tests backend   & \codigo{pytest tests/}        & \ok\ 235 pasan, cob. 83.38\% \\
    Lint backend    & \codigo{ruff check app/ tests/} & \ok\ limpio \\
    Tipos frontend  & \codigo{vue-tsc --noEmit}      & \ok\ limpio \\
    Lint frontend   & \codigo{npm run lint}          & \ok\ 0 errores, 17 advertencias \\
    Tests frontend  & \codigo{vitest --run}          & \ok\ 265 pasan, 9 omitidos$^{\ast}$ \\
    Tipos backend   & \codigo{mypy app/}             & \ok\ limpio (33 archivos) \\
    \bottomrule
  \end{tabular}
  \caption{Panel de calidad. $^{\ast}$Vitest 4; el ajuste de
    \codigo{localStorage} de \codigo{src/test/setup.ts} ya forma parte
    del árbol, sin workarounds de entorno.}
  \label{tab:calidad}
\end{table}

\section{Puertas de calidad configuradas}
\begin{itemize}
  \item Cobertura backend $\ge 80\%$ (\codigo{--cov-fail-under},
        \codigo{pyproject.toml}); cobertura actual 83.38\% en corrida
        limpia (un comentario del propio archivo menciona un objetivo
        de 95\% de una fase previa del plan maestro).
  \item Umbrales de cobertura frontend en \codigo{vitest.config.ts}
        (líneas/sentencias 95\%, ramas 83\%, funciones 70\%).
  \item Guardia de deriva OpenAPI: \codigo{openapi.json} se regenera y
        se compara en CI.
  \item Auditoría de design tokens:
        \codigo{scripts/audit_design_tokens.sh} limita los valores en
        duro fuera de \codigo{tokens.css}.
  \item Escaneo de vulnerabilidades con Trivy (imágenes y filesystem)
        en el job de Docker.
  \item Humo post-despliegue:
        \codigo{backend/scripts/smoke.sh} (código de salida no
        cero al primer fallo) y verificador extremo a extremo
        \codigo{scripts/verify_api.py} que cruza cada endpoint contra
        la base en disco.
  \item Pre-commit: ruff (lint+formato), mypy, ganchos estándar.
  \item Réplica local del CI: \codigo{scripts/run_ci_local.sh}.
\end{itemize}

\section{Jobs de CI (GitHub Actions)}
\codigo{lint}, \codigo{typecheck} (build con vue-tsc), \codigo{test},
\codigo{fastapi-checks} (ruff + pytest + deriva OpenAPI) y
\codigo{docker-build} / \codigo{docker-build-fastapi} (build + health +
smoke + Trivy + verificador E2E).

% =====================================================================
\chapter{Estado del proyecto y trabajo pendiente}
\label{cap:pendientes}
% =====================================================================

\section{Veredicto}
El sistema es \textbf{funcionalmente completo, está desplegado} en
producción (\codigo{cenidim.\allowbreak darylemb.dev}) y las suites de
backend y frontend están íntegramente en verde. Los bloqueantes técnicos
detectados en la revisión fueron resueltos en la rama \codigo{release}
(\S\ref{sec:resueltos}); lo que queda es trabajo menor y las decisiones
editoriales del catálogo.

\section{Resueltos en la rama \codigo{release}}
\label{sec:resueltos}
\begin{itemize}
  \item \textbf{Tipado backend.} \codigo{mypy app/} reportaba 41
        errores en 7 archivos. Fueron corregidos con anotaciones reales
        (no supresiones): generadores asíncronos tipados
        (\codigo{session\_scope}, \codigo{get\_db}), helpers SQL de
        \codigo{public.py} con \codigo{TypeVar} acotado a
        \codigo{Select}, \codigo{cast} a \codigo{CursorResult} para
        \codigo{lastrowid}/\codigo{rowcount}, y firma anotada del
        middleware CSRF. Resultado: \ok\ limpio en 33 archivos, y el
        hook de pre-commit ya no bloquea.
  \item \textbf{Tests frontend en Node $\ge$ 24.} El
        \codigo{localStorage} nativo de Node sombreaba al de jsdom (63
        fallos). \codigo{src/test/setup.ts} instala ahora un polyfill
        \codigo{Storage} en memoria; los 265 tests pasan sin
        \codigo{NODE\_OPTIONS} ni workarounds de entorno.
  \item \textbf{Dependencias frontend.} \codigo{npm audit} reportaba 15
        vulnerabilidades (1 crítica). Se eliminó \codigo{vue-wordcloud}
        (sin uso: la nube es el componente propio
        \codigo{WordCloud.vue}), se actualizó vitest a la 4.1.11
        (corrigiendo el mock reset en \codigo{auth.spec.ts} conforme al
        nuevo comportamiento de \codigo{restoreAllMocks}) y se aplicó
        \codigo{npm audit fix}. Resultado: \ok\ 0 vulnerabilidades.
  \item \textbf{Higiene del repositorio.} Los CSV de clasificación
        manual se conservaron en \codigo{historical/} con un README de
        procedencia (no son entrada del pipeline).
        \codigo{frontend/\allowbreak dist/} dejó de trackearse (lo
        regenera el Dockerfile desde fuente). Las dos
        \codigo{letras.db} se mantienen trackeadas de forma deliberada,
        como referencia para nuevos mantenedores. \codigo{.gitignore}
        corregido.
  \item \textbf{CI en \codigo{release}.} El workflow de GitHub Actions
        ahora dispara también en la rama \codigo{release} (push y PR).
\end{itemize}

\section{Pendientes técnicos}
\begin{longtable}{p{0.9cm}p{3.4cm}p{9.4cm}}
  \toprule
  & \textbf{Área} & \textbf{Pendiente} \\
  \midrule
  \endhead
  \pend & Licencia &
    Sigue sin existir un archivo \codigo{LICENSE}; para la entrega
    formal conviene elegir una (o confirmar ``todos los derechos
    reservados'' con el cliente). \\
  \pend & Observabilidad &
    Deprecación FastAPI: \codigo{on_event} en \codigo{app/main.py}
    debe migrarse a manejadores \emph{lifespan} (348 avisos en la
    corrida de pytest). \\
  \pend & Configuración &
    El comentario de cobertura en \codigo{pyproject.toml} cita un
    objetivo de 95\% de una fase previa, pero la puerta vigente es
    80\% (cobertura actual 83.38\%); conviene alinearlos. \\
  \bottomrule
  \caption{Pendientes técnicos restantes tras la rama \codigo{release}.}
  \label{tab:pendientes-tec}
\end{longtable}

\section{Pendientes editoriales (con el equipo del archivo)}
\begin{itemize}
  \item Catálogo canónico de temas: colapsar conceptos distintos con
        nombres parecidos (los meros errores tipográficos ya se pliegan
        con \codigo{TEMA_TYPO_MAP}).
  \item Homogeneizar la terminología del tablero con los documentos
        de trabajo y el ensayo (las etiquetas \emph{in-app} ya tienen
        explicaciones \raisebox{0pt}{$\circledast$}\,(i)).
  \item \begin{sloppypar}Revisar el listado de canciones sin catalogar
        (documento \codigo{CANCIONES-SIN-CATALOGAR-2026-08-03.md} en
        \codigo{docs/}, retirado del árbol en \codigo{ff4303b} y
        recuperable de la historia de git) para el acervo aún sin
        tema.\end{sloppypar}
\end{itemize}

\section{Riesgos y notas de operación}
\begin{itemize}
  \item \codigo{CENIDIM_JWT_SECRET} es obligatorio en producción; el
        backend no arranca sin él.
  \item \codigo{letras.db} es regenerable: ante corrupción basta
        reiniciar el stack (\codigo{db-init} lo reconstruye).
  \item El rollback a Go ya no existe como árbol; solo queda el commit
        congelado \codigo{2aab765} en la historia.
\end{itemize}

% =====================================================================
\appendix
\chapter{Glosario breve}
% =====================================================================

\begin{description}[style=nextline]
  \item[Fonograma] Registro de álbum/disco fuente del que provienen las
    canciones; unidad de colección (280 en el corpus).
  \item[OOV] \emph{Out Of Vocabulary}: porcentaje de palabras de una
    letra que el modelo spaCy no reconoce. Proxy de qué tan cotidiano
    o especializado es el lenguaje.
  \item[Tema] Descriptor editorial asignado por catalogadores humanos
    (línea \codigo{Tema:} del archivo fuente). Nunca se infiere por
    palabras clave.
  \item[Clasificación lingüística] Tipología automática por letra:
    español estándar, español regional o lengua indígena (umbrales OOV
    en \S\ref{cap:pipeline}).
  \item[db-init] Servicio sidecar de Docker que reconstruye
    \codigo{letras.db} (build $\to$ clasifica $\to$ normaliza) antes de
    que arranque el backend.
  \item[email\_outbox] Tabla-bandeja donde caen los correos de
    recuperación cuando no hay \codigo{RESEND_API_KEY} (desarrollo).
  \item[ADR] \emph{Architecture Decision Record}; el de este proyecto
    es \codigo{docs/adr/\allowbreak 0001-\allowbreak fastapi-\allowbreak replaces-\allowbreak go.md}
    (retirado del árbol en \codigo{ff4303b}; en la historia de git).
\end{description}

El glosario canónico completo, con la definición de cada KPI y gráfica
del tablero, era \codigo{docs/GLOSARIO.md} (retirado en
\codigo{ff4303b}); en la aplicación, la fuente viva es
\codigo{frontend/src/\allowbreak config/\allowbreak chartInfo.ts}.

% =====================================================================
\begin{thebibliography}{9}
\addcontentsline{toc}{chapter}{Referencias}
% =====================================================================

\bibitem{readme} Repositorio del proyecto,
  \textit{README.md} (raíz) y \textit{backend/README.md}.

\bibitem{adr} ADR 0001: \textit{FastAPI reemplaza a Go},
  \codigo{docs/adr/0001-fastapi-replaces-go.md} (retirado del árbol en
  \codigo{ff4303b}; recuperable de la historia de git).

\bibitem{cutover} \textit{Cutover Playbook --- Fases 7 y 9},
  \codigo{docs/CUTOVER.md} (ídem).

\bibitem{glosario} \textit{Glosario del tablero},
  \codigo{docs/GLOSARIO.md} (ídem; la fuente viva en la aplicación es
  \codigo{frontend/src/config/chartInfo.ts}).

\bibitem{review} \textit{Respuestas a las 16 observaciones del tablero},
  \codigo{docs/REVIEW-2026-07-01-STATUS.md} y
  \codigo{docs/RESPUESTAS-CLIENTE-CENIDIM-2026-07-28.md} (ídem).

\bibitem{openapi} Especificación OpenAPI 3.1 de la API,
  \codigo{backend/openapi.json}.

\bibitem{fastapi} Documentación de FastAPI,
  \url{https://fastapi.tiangolo.com/}.

\bibitem{vue} Documentación de Vue 3, \url{https://vuejs.org/};
  Pinia, \url{https://pinia.vuejs.org/};
  Vue Router, \url{https://router.vuejs.org/}.

\bibitem{spacy} spaCy, modelo español \codigo{es_core_news_md},
  \url{https://spacy.io/models/es}.

\bibitem{chartjs} Chart.js, \url{https://www.chartjs.org/};
  envoltorio \codigo{vue-chartjs}.

\end{thebibliography}

\end{document}
